Content provenance standards attach a signed record of origin and edits to a media file. The C2PA specification \cite{c2pa2026spec} binds that record to the file cryptographically, so the binding holds only as long as the bytes do. Re-encoding, resizing, a screenshot or a platform's upload pipeline strips or invalidates the manifest, and the content then circulates without its history. C2PA therefore defines \emph{soft bindings}: identifiers derived from the content that let anyone holding a copy look the manifest up again. A soft binding is either an invisible \emph{watermark}, which carries a key inside the pixels or samples, or a \emph{fingerprint}, which is computed from the content and matched against a registry. Regulation adds a second reason to put a signal into the content itself: Article~50(2) of the EU Artificial Intelligence Act requires providers of generative systems to mark synthetic audio, images, video and text in a machine-readable, detectable form, with solutions that are effective, interoperable, robust and reliable as far as technically feasible \cite{eu2024aiact, eu2026omnibus}. Neither the standard nor the regulation prescribes a method or a performance threshold.

The two families are complementary. A watermark answers directly and needs no index, but it must be embedded before release, it can be removed, and it says nothing about content that was never marked. A fingerprint needs no change to the content and applies to an existing archive, but every lookup risks binding a query to the wrong asset, and that risk grows with the registry: a service holding millions of assets needs false-match rates per comparison of $10^{-7}$ or lower. For both families, the rate at which unmarked or unregistered content is reported as a match matters as much as robustness, because a false match attaches someone else's provenance to a piece of content. That rate depends on the decision rule a verifier applies, not only on the model.

Choosing either kind of soft binding is therefore an empirical question, and the published evidence is fragmented. Each new watermark is reported against baselines of its authors' choosing, at its own embedding strength and with its own metrics \cite{bui2025trustmark, xu2025invismark, soucek2025pixelseal, sanroman2024audioseal}; watermark benchmarks cover one modality each \cite{an2024waves, lu2025vine, liu2024audiomarkbench, jiang2025videomarkbench}. Copy-detection benchmarks rank learned descriptors by average precision without a fixed false-match rate \cite{douze2021isc}, perceptual-hash studies measure robustness to benign edits or, separately, to adversaries \cite{jain2022adversarial, prokos2023squint}, and the one open fingerprint on the C2PA soft-binding list, the ISCC of ISO 24138 \cite{iso24138}, has not been compared with either. Licences decide the matter for many adopters and differ between a paper and its release: the strongest learned image descriptors ship with weights trained on non-commercial data, and some watermark checkpoints carry no licence at all. Finally, the two families are evaluated separately, although a provenance service would deploy them together.

This paper benchmarks openly available watermarks and fingerprints for images, audio and video under one protocol, designed around the decisions a provenance service must make. Our contributions are as follows.
\begin{itemize}
\item A watermarking track of \val{wm:count/image} image, \val{wm:count/audio} audio and \val{wm:count/video} video configurations from \val{wm:count/families} methods on public, licence-audited media, measuring perceptual quality with display-aware metrics, robustness to \val{wm:count/transforms} routine transformations with worst-subset reporting, payload recovery, false positives under blind and expected-payload decision rules, and embedding latency and estimated cloud cost (Sections~\ref{sec:wm-method} and~\ref{sec:wm-results}).
\item A fingerprinting track of \val{fp:count/image/methods} image, \val{fp:count/audio/methods} audio and \val{fp:count/video/methods} video methods with \val{fp:count/attacks} transformations, a registry of about a hundred thousand images, and thresholds fixed on held-out negatives at query-level and registry-size-independent pair-level false-match rates down to $10^{-7}$; a partial-edit track with localization; a test of second-stage geometric verification; and white-box, transfer, forgery and inversion attacks (Sections~\ref{sec:fp-method} and~\ref{sec:fp-results}).
\item A measurement of how the two families complement each other on identical media and attacks, and of how embedding a watermark moves fingerprints (Section~\ref{sec:together}).
\item Statistics that treat the source item, not the transformed copy, as the independent unit, and two reproducibility findings: detector outputs that change between x86 and ARM hosts because colour conversion differs, and benchmark inputs that must be pinned beyond the corpus itself.
\end{itemize}
Section~\ref{sec:related} reviews related work, Section~\ref{sec:design} describes the design shared by both tracks and each track's protocol, Sections~\ref{sec:wm-results}--\ref{sec:together} report the results, and Section~\ref{sec:discussion} discusses what they imply for a C2PA soft-binding service.
